\chapter{Conclusions and Outlook}
\label{ch:conclusions-and-outlook}

\textbf{Status: synthesis of the present evidence boundary and proposed work.
No completed production reduction or physical validation is claimed.}

The central question of this monograph is whether a first-principles
Kohn--Sham description can be replaced by a smaller model while preserving the
information required for a declared use. The current project state does not yet
supply a numerical answer. It does, however, establish the distinctions,
contracts, and evidence boundaries needed for that answer to be meaningful.

\section{Present conclusions}

Several conclusions are supported at the present evidence level.

First, spectral fitting and represented-operator fitting are different
requirements. Eigenvalues do not identify localized orbitals, gauge, site and
orbital correspondence, or the matrix elements through which later impurity
operators act. A reduced model intended only for spectral interpolation may be
adequate for that use while remaining unsuitable as a common coordinate system
for perturbations.

Second, pristine and doped matrix subtraction is not defined by equal array
shape. The retained state spaces, basis ordering, geometry, units, spin
convention, energy zero, and gauge must agree or be related by an explicit
identification map. A difference formed before these prerequisites are met has
no established impurity-operator meaning.

Third, projection, disentanglement, gauge transformation, Wannier localization,
and real-space truncation are separate operations. Their assumptions and losses
must remain attributable. In particular, localization does not itself establish
a low-rank approximation or authorize an unstated hopping cutoff.

Fourth, parent-model, numerical, and reduction errors answer different
questions. The PBE parent's relationship to experiment is not the same as a
finite-cutoff error, and neither is the same as the error introduced by an
$sp^3s^{\ast}$ model class. Combining them without a compatible mathematical
contract would hide rather than explain the quality of the reduction.

Finally, the repository has an accepted represented-operator software
foundation and retained tutorial and bootstrap execution observations. Those
artifacts support their bounded software and workflow claims. They do not yet
supply the converged production parent, Wannier operator, reduced bulk model,
impurity operator, or continuum result needed to answer the scientific
questions.

\section{The unresolved bulk question}

The first unresolved scientific question is whether one member of a prescribed
bulk tight-binding hierarchy can satisfy both spectral and aligned-operator
criteria relative to the same accepted Kohn--Sham parent. The answer may take
one of three forms:

\begin{enumerate}
  \item the first tested class contains a common admissible Hamiltonian;
  \item a later, more expressive class is the first to contain one; or
  \item no tested class satisfies the joint criteria.
\end{enumerate}

None of these outcomes is assumed. If no common model is found, the set
separation and its orbital-, symmetry-, or shell-resolved contributions may be
more informative than forcing an expanded fit. Such a result would identify the
information that the tested class cannot preserve.

\section{Near-term bulk program}

The next scientific sequence remains within pristine silicon:

\begin{enumerate}
  \item complete the canonical convergence boundary for the declared
    PBE/PseudoDojo parent;
  \item determine and accept the internally consistent equilibrium lattice
    reference;
  \item construct the production SCF parent and purpose-specific band-path,
    valley, and mass children;
  \item freeze the target subspace, QE--Wannier interface, and validated Wannier
    operator;
  \item verify the alignment machinery first on controlled transformations and
    then on independent bulk constructions; and
  \item compare direct and Wannier-mediated reductions on frozen training and
    withheld datasets from the same parent.
\end{enumerate}

This is a scientific dependency sequence, not execution authorization. Each
protected calculation still requires its applicable activation, preflight,
resource report, and human authorization.

\section{Phosphorus as the first impurity demonstration}

After the bulk representation and alignment gates pass, phosphorus is the first
complete impurity path. The branch requires a compatible P pseudopotential,
neutral spin-polarized P$_\mathrm{Si}^{0}$ supercells, controlled internal
relaxation, matched pristine comparison cells, and supercell-size evidence. The
aligned difference is then decomposed into nested scalar, orbital, onsite,
hopping, and range-restricted model classes.

The objective is not to assume that a screened Coulomb potential plus one
central-cell parameter is sufficient. It is to determine which operator and
bound-state information each tested class preserves. A continuum donor model
becomes a first-principles reduction only after its host parameters, embedding,
exterior error, cross-coupling error, and bound-state comparison are tied to the
accepted atomistic parent.

\section{Boron as a transferability branch}

Boron tests whether the methods transfer to a physically different band-edge
problem. The final acceptor branch requires a fully relativistic noncollinear
spinor treatment with SOC and a compatible valence subspace. The phosphorus
basis or scalar model class cannot be transferred unchanged merely for
convenience.

Comparing the P and B hierarchies may reveal which reduction features are common
to substitutional impurities and which are specific to conduction- or
valence-manifold structure. Such a comparison becomes eligible only after each
branch has its own accepted atomistic operator and minimal-model evidence.

\section{Composition and path dependence}

A completed hierarchy still requires cross-path assessment. Direct and
Wannier-mediated bulk fits may disagree. Extracting an impurity before reduction
may differ from reducing pristine and doped parents before subtraction.
Successive lattice and continuum reductions may amplify earlier errors.

These noncommuting diagrams are not expected to vanish by definition. Their
paths, aligned outputs, norms, and validation domains must be measured. The
final compositional claim may be that a diagram commutes within tolerance over a
declared domain, that it fails by a quantified defect, or that available
metadata are insufficient to compare the paths.

\section{Negative results as supported outcomes}

The research program is designed to retain negative findings. An unacceptable
nearest-neighbor model, an unstable alignment, a nonconverged impurity sequence,
an empty admissible-set intersection, or an infinite crossover radius over the
tested domain can all be valid outcomes. Their value depends on preserving the
frozen model hierarchy, rejected candidates, tested range, metrics, and
limitations.

A negative finding is not generalized beyond its tested classes and domain. A
failed software contract is corrected; a failed numerical tolerance is refined
or reported; a physically inadequate model class is not repaired by weakening
its acceptance rule after the result is known.

\section{Scholarly outputs}

The monograph remains the broad account of definitions, methods, provenance,
results, and limitations. Article manuscripts extract narrower questions.
P01, for example, can focus on bulk model-class expressiveness and the
compatibility of spectral and operator criteria without claiming completion of
the impurity program. Later articles may address alignment, impurity-operator
hierarchies, or continuum crossover only when their evidence exists.

Presentations may use the same retained results at a different explanatory
level, but they do not change evidentiary status. A conference figure derived
from tutorial data remains tutorial evidence; a proposed workflow diagram
remains proposed work; a production result retains its calculation identity and
limitations.

\section{Final perspective}

The intended contribution is not a single fitted parameter table. It is a
controlled account of when representations can be identified, when reduced
operators preserve their declared parent information, and where approximation
paths cease to agree. Reaching that account requires the full chain from
physical specification through numerical convergence, representation,
alignment, reduction, and evidence.

At present, the project has established the framework and bounded software
foundation but has not completed the scientific reduction sequence. The
remaining work is therefore stated as work: construct the accepted parent,
measure each loss, retain negative outcomes, and draw conclusions only at the
strength supported by the resulting evidence.
